% Part: normal-modal-logic
% Chapter: syntax-and-semantics
% Section: introduction

\documentclass[../../../include/open-logic-section]{subfiles}

\begin{document}

\olfileid{nml}{syn}{int}

\olsection{پېژندنه}

مودالي منطق له \emph{مودالي قضيو} او د هغو تر منځ د استلزام له
اړيکو سره سروکار لري۔ د مودالي قضيو بېلګې دا دي:
\begin{enumerate}
\item دا لازمه ده چې $2+2=4$.
\item دا لازمه ده چې سبا د باران کېدو امکان وي.
\item که دا لازمه وي چې~$!A$ ممکن وي، نو دا ممکنه ده چې~$!A$.
\end{enumerate}
امکان او لزوم يوازيني موداليتونه نۀ دي: نور يوځايي نښلوونکي هم
موداليتونه ګڼل کېږي؛ لکه «پکار ده چې~$!A$ رښتيا وي»، «داسې به وي
چې~$!A$»، «ډېنا پوهېږي چې~$!A$»، يا «ډېنا باور لري چې~$!A$».

مودالي منطق لومړے ځل د ارسطو په \emph{De Interpretatione} اثر کښې
راڅرګند شو: هغه لومړے کس و چې پام ئې وکړ لزوم د امکان موجب کېږي، خو
برعکس نۀ؛ امکان او لزوم يو د بل په وسيله تعريفېدلے شي؛ که
$!A \land !B$ په ممکنه توګه رښتيا وي، نو $!A$ او $!B$ هر يو په ممکنه
توګه رښتيا دي، خو برعکس نۀ؛ او که $!A \to !B$ لازم وي، نو که
$!A$ لازم وي، $!B$ هم لازم وي.

مودالي منطق ته لومړنۍ عصري لاره د سي.~آی. لېويس کار و، چې د لېويس
او لنګفورډ په \emph{سمبوليک منطق} (1932) کښې اوج ته ورسېد۔ لېويس او
لنګفورډ د مادي شرطيې په وسيله د استلزام له ښودنې راضي نۀ وو:
$!A \lif !B$ د «له $!A$ څخه $!B$ لازمېږي» کمزورے بديل دے۔ د هغې
پر ځاے ئې استلزام داسې مشخص کړ: «ضرورتاً، که $!A$ وي نو $!B$ وي»؛
نښه ئې $!A \strictif !B$ ده۔ د بېلابېلو خاصيتونو د روښانولو په هڅه
کښې لېويس پنځه بېل مودالي نظامونه، \Log{S1}, \ldots, \Log{S4},
\Log{S5}، وپېژندل، چې وروستي دوه لا هم کارېږي.

د لېويس او لنګفورډ لاره خالصه \emph{نحوي} وه: معقول بديهي اصول او
قاعدې ئې وټاکلې او وئې څېړل چې د هغو په مټ څه ثابتېدلے شي۔ معنايي
لار اوږده موده لاسرسۍ ته نۀ وه، تر هغه چې رودولف کارناپ په
\emph{معنٰی او لزوم} (1947) کښې \emph{د حالت توصيف} مفهوم وکاراوه،
يعنې د اټومي جملو يوه ټولګه (هغه جملې چې په هماغه توصيف کښې
«رښتيا» دي)۔ کارناپ چې د صدق تعريف هرې جملې $!A$ ته وغځاوه، نو
$!A$ هله \emph{لازماً رښتيا} بولي چې د حالت په ټولو توصيفونو کښې
رښتيا وي۔ د کارناپ لاره \emph{پرله پسې} موداليتونه نۀ شي سمبالولے،
ځکه چې د «ممکن، ضرورتاً، \ldots ممکن $!A$» په بڼه جملې تل تر ټولو
دننني موداليت ته راټولېږي.

په مودالي معناپوهنه کښې ستر پرمختګ د ساول کرېپکي له مقالې «په
مودالي منطق کښې د بشپړتيا يوه قضيه» (JSL، 1959) سره راغے۔ کرېپکي
خپل کار د لايبنېڅ پر دې نظر ولاړ کړ چې يو بيان هله لازماً رښتيا وي
چې «په ټولو ممکنه نړۍو کښې» رښتيا وي۔ خو دا نظر د کارناپ د لارې
هماغه نيمګړتيا لري: په يوه نړۍ $w$ (يا د حالت په توصيف $s$) کښې د
بيان صدق په خپله له $w$ سره تړاو نۀ لري۔ نو کرېپکي ومانله چې
نړۍوې د \emph{لاسرسي اړيکې} $R$ له لارې يو له بل سره تړلې دي؛ او
«ضرورتاً $!A$» په نړۍ $w$ کښې هله او يوازې هله رښتيا وي چې $!A$
په ټولو هغو نړۍو $w'$ کښې رښتيا وي چې له $w$ څخه \emph{ورته لاسرسے
شته}۔ هغه معناپوهنې چې د دې لارې يوه بڼه کاروي، د کرېپکي معناپوهنې
بلل کېږي او د مودالي منطقونو (په جمع بڼه) د پراخې ودې لار ئې پرانيسته.

کله چې مودالي منطق د کرېپکي د معناپوهنې له مخې تفسير شي، راته ښيي
چې \emph{اړيکيز جوړښتونه} له «دننه» څنګه ښکاري۔ اړيکيز جوړښت يوازې
يو سټ دے چې يوه دوه ځايه اړيکه ورباندې ټاکل شوې وي (د بېلګې په
توګه، په ټولګي کښې د زده کوونکو هغه سټ چې د ټولنيز امنيت د شمېرو له
مخې مرتب شوے وي، يو اړيکيز جوړښت دے)۔ خو اړيکيز جوړښتونه په
بېلابېلو ډګرونو کښې راځي: د نړۍ د حالتونو د نسبي امکان تر څنګ، د
يوه عامل معرفتي حالتونه د معرفتي امکان له مخې تړلے کېدلے شي، يا د
يوه خوځنده نظام حالتونه د هغو تر منځ د لېږدونو له مخې، او داسې نور۔
مودالي منطق د دې ټولو د مدل جوړولو دپاره کارېږي: لومړے عادي، د
حقيقت اړوند مودالي منطق راکوي؛ نور معرفتي منطق، خوځنده منطق، او
داسې نور.

موږ پر يوه ځانګړي اړخ تمرکز کوو چې مودالي منطقپوهان ئې «د مطابقت
تيوري» بولي۔ د کرېپکي له لومړنيو مهمو موندنو څخه يوه دا وه چې د
لاسرسي اړيکې~$R$ ډېر خاصيتونه (لکه تعدي، تقارن او داسې نور) پخپله
\emph{په مودالي ژبه کښې} د مناسبو «مودالي طرحو» په وسيله مشخصېدلے
شي۔ د بېلګې په توګه، مودالي منطقپوهان وايي چې د $R$ انعکاسي والی له
دې طرحې سره «مطابقت لري»: «که $!A$ لازماً رښتيا وي، نو $!A$ رښتيا
وي»۔ موږ په عمده توګه د مودالي منطق د هغو کلاسيکي نظامونو (لکه
\Log{S4} او \Log{S5}) د مطابقت تيوري څېړو چې د \Ax{D}, \Ax{T},
\Ax{B}, \Ax{4} او~\Ax{5} طرحو په ګډولو سره ترلاسه کېږي.

\end{document}
